\section{Introduction}
\label{sec:intro}
% Diffusion-based language models (DLLMs) offer a promising alternative to autoregressive (AR) generation by predicting multiple tokens in parallel through iterative denoising. However, existing DLLMs suffer from three critical gaps that prevent practical adoption: (1)~training from scratch is prohibitively expensive in both resource and data,  and yields models that lag behind AR counterparts in quality; (2)~inference relies on confidence-based thresholds that lack formal quality guarantees; and (3)~the iterative denoising nature of DLLMs is incompatible with existing AR serving infrastructure, forcing practitioners to build custom systems that cannot leverage the mature, well-optimized AR inference stacks.

Autoregressive (AR) models have become the dominant paradigm for large language models, achieving remarkable quality through next-token prediction at scale~\citep{qwen3, deepseekv3}. However, their strictly sequential decoding creates an inherent throughput bottleneck: each token must wait for all preceding tokens, leaving modern parallel hardware underutilized during generation. Diffusion-based language models (DLLMs)~\citep{austin2021d3pm, sahoo2024mdlm, nie2025llada} offer an appealing alternative by predicting multiple tokens in parallel through iterative denoising. Yet despite growing research interest, DLLMs have not been widely adopted in practice.

\yifan{Revised per Fan's comments: reordered bottlenecks (quality first, systems last); defined prefill-decode mismatch.}
We identify four fundamental bottlenecks. First, DLLM training extracts far less supervision signal per FLOP than AR training due to masking, resulting in a persistent accuracy gap even with conversion-based approaches. Second, DLLMs suffer from a prefill--decoding gap---prefill computes the KV cache but provides no direct prediction signal for the first generated token, unlike AR where the last prefill hidden state naturally initiates generation. Third, parallel decoding purchases latency reduction at the cost of significantly more total compute, making DLLMs less cost-effective than AR models under high concurrency. Fourth, the entire modern serving stack---attention kernels, KV cache management, batching, and pipelining---is built around AR decoding; the denoising loop of DLLMs does not fit naturally into these systems. We analyze these bottlenecks quantitatively in Section~\ref{sec:motivation}.

Rather than adapting the ecosystem to diffusion---as prior efforts such as FastDLLM~\citep{fastdllm2025} and WeDLM~\citep{wedlm2025} attempt---we take the opposite perspective: we \emph{start from the AR side} and ask what useful ideas can be borrowed from diffusion to make AR decoding more parallel. This viewpoint is shared by multi-token prediction~\citep{gloeckle2024mtp, deepseekv3}, speculative decoding~\citep{leviathan2023speculative, eagle3}, and register-based AR variants~\citep{cllms2024}, which also inject parallelism into AR generation.
Prior work on AR-to-diffusion conversion has moved in this direction: SDAR~\citep{sdar2025} demonstrated that converting pretrained AR models into block-diffusion generators is far more efficient than training from scratch, and NBDiff~\citep{nbdiff2025} went further by adopting context-causal attention masks to better preserve the AR structure. 

\yifan{Revised per Fan's comment: strengthened contribution claims with SOTA quality, baseline comparisons, latency/throughput numbers, and industrial-scale concurrency.}
We argue that the key to unlocking the full potential of AR models as diffusion LLMs is to \emph{sufficiently respect the knowledge and patterns that AR pretraining has fostered}---in attention structure, training objective, and serving infrastructure alike. To this end, we introduce \textbf{Strided DLLM}, a flexible diffusion language model that generates $N$ tokens per step---analogous to strided convolution, where the generation stride controls the parallelism--quality trade-off. Strided DLLM achieves state-of-the-art quality among diffusion LLMs: to the best of our knowledge, it is the first to achieve comparable quality to its AR base model while delivering \textbf{XX}$\times$ speedup in latency and \textbf{XX}$\times$ improvement in throughput. With only 8B parameters, it surpasses LLaDA-2.1-mini (16B) across all benchmarks and outperforms competitive speculative decoding baselines such as EAGLE-3 and DFlash in both latency and throughput under high concurrency. Critically, unlike prior DLLMs whose throughput collapses at scale, Strided DLLM sustains \textbf{XX} tok/s at batch size 64---\textbf{XX}\% higher than the AR baseline---making it viable for industrial-scale deployment. Concretely, our contributions are:

\begin{denseitemize}
    \item \textbf{Strided DLLM: a flexible diffusion language model.} We propose a novel method to convert pretrained AR models (Qwen3-8B and Qwen3-32B) into strided diffusion models that can decode at variable block sizes. Logit-shifted causal training introduces minimal divergence from the pretrained AR objective, so the model remains robust across stride configurations---including the all-masked setting. This removes the need for distillation schedules or masking curricula, yielding a data-efficient conversion that preserves the base model's knowledge.

    \item \textbf{Introspective Strided Decoding (ISD).} Rather than relying on confidence-based thresholds, we propose ISD, where the model introspects on its own parallel proposals by verifying them against the causal (AR) distribution via a $p/q$ acceptance criterion. This provides a mathematical guarantee that the output matches the base AR model's distribution---without confidence heuristics or an external draft model.

    \item \textbf{Lossless ISD with residual adaptation.} LoRA adapters are active only on MASK positions during strided decoding; verification positions use base-model-only weights. This yields bit-for-bit identical outputs to the base AR model---a true lossless acceleration. The approach supports both near-lossless (merged LoRA) and true lossless (gated LoRA) modes.

    \item \textbf{AR-compatible serving infrastructure.} Our design maximally reuses existing AR infrastructure (vLLM/SGLang) through causal masking and fused Triton kernels, rather than building a custom DLLM serving stack. The system achieves higher tokens-per-second than all baselines, including EAGLE-3 and LLaDA-2.1-mini.
\end{denseitemize}

\yifan{Added per Fan's comment.}
We will open-source model weights, training code, and SGLang integration to facilitate community research and deployment.